DDPM trains a noise-prediction model and generates samples through a learned
reverse process \cite{Ho2020}. Ho et al.\ report parameter EMA with a fixed
decay without a decay sweep in their experimental details. Nichol and
Dhariwal introduce a cosine schedule for the forward diffusion noise
\cite{Nichol2021}. That schedule concerns corruption over diffusion time. It
is distinct from the optimizer learning-rate policy over training updates,
which is the intervention in the retained experiment.

Stochastic Weight Averaging (SWA) combines selected SGD iterates under constant
or cyclical learning rates \cite{Izmailov2018}. Its averaging rule differs
from a fixed exponential decay applied after every AdamW update. SWA is a
relevant alternative that this experiment did not test, so the results cannot
rank EMA against it.

Morales-Brotons et al.\ study EMA dynamics and report strong intermediate
performance before full learning-rate reduction in their classification
experiments \cite{MoralesBrotons2024}. Their validation-based early stopping,
SGD trajectories and batch-normalization choices differ from the fixed
endpoints here. Ajroldi et al.\ subsequently compare averaging across a wider
optimization benchmark and examine no, partial and full annealing
\cite{Ajroldi2025}. Their results show that averaging gains can diminish as
the baseline is fully annealed and that combining the two can give the best
final performance in the settings tested. This is direct prior evidence for
the qualitative relationship considered in this paper.

For diffusion models, Karras et al.\ show that preferred averaging profiles
depend on architecture and training time \cite{Karras2024}. Their
power-function and post-hoc averaging procedures, image-scale models and FID
evaluation differ from the two fixed decays and SW1 metric here. Their findings
support caution about transferring a preferred averaging window across designs.

The literature also contains contrary ordinary-EMA outcomes. Li et al.\ report
video-prediction and regression cells where ordinary EMA performs worse than
the basic model, while proposing a switching procedure that feeds averaged
parameters back into training \cite{Li2024}. That intervention changes the
training trajectory. Its reported gains and the unhelpful bootstrapping
experiment in Morales-Brotons et al.\ are not matched comparisons. The present
study evaluates averages of the same raw trajectories and cannot resolve that
protocol-dependent difference.
